\chapter{Automated Options Market Making}\label{ch:hf:automated-options-market-making}

An options market maker quotes every strike and expiry of a stock at once, several thousand prices on a dozen exchanges, and a sweep can hit forty
of them in the same millisecond; the exchange's protection is what stops the forty-first. In this chapter's chain of 170 series, a trader who
knows the stock has moved 1\% and hits every stale quote, one every five microseconds, would take \$79\,000 from a market maker without protection;
a protection that pulls all quotes after thirteen executions holds the loss to \$12\,900 and, on an ordinary day of multi-leg customer orders, costs
three fills out of 8\,235.

\section{The live surface}

An options market maker's theoretical values come from a volatility surface (One Quant Book 5, chapters 7 and 8): a parametrised smile per
expiry, fitted to the market's quotes and trades, moved with the underlying according to a rule (sticky strike, sticky delta) and refitted when
the market says the rule was wrong.

\begin{definition}[Live surface fit]\label{def:hf:automated-options-market-making:live}
A \emph{live surface fit}\index{live surface fit} is the continuous re-estimation of an options market maker's volatility surface during the trading
day, updating each expiry's parameters from the latest quotes and trades as they arrive, fast enough that every quote the market maker sends is
priced on a surface consistent with the current market.
\end{definition}

At high frequency the work is split: the underlying's price moves the whole surface through the rule, microsecond by microsecond, and the smile's
shape is refitted on a slower clock, typically when enough new information has arrived. A raw SVI slice fitted by the quasi-explicit method
(Book 5's \texttt{firm.svi}) to seventeen quoted volatilities, each with 0.2 volatility points of noise, recovers the true smile to 0.025 points: the
fit averages the noise that individual quotes carry. Muravyev and Pearson showed why the fair value matters so much to takers too: option prices
change predictably at high frequency, many traders time their executions to buy when the fair value is close to the ask and sell when it is close to
the bid, and the effective spreads of those who time are under 40\% of conventional measures.

\section{Mass quoting and quote protection}

\begin{definition}[Mass quote]\label{def:hf:automated-options-market-making:mass}
A \emph{mass quote}\index{mass quote} is a single message in which a market maker sends or replaces bids and offers in many option series at
once, which the exchange processes as a unit, so that a market maker can update thousands of quotes after a move in the underlying without sending
thousands of separate orders.
\end{definition}

A chain of 17 strikes and five expiries, calls and puts, is 170 series; quoted at half a volatility point either side of theo
(\cref{lst:hf:automated-options-market-making:quote}), the 30-day at-the-money call at \$2.92 is 11 cents wide, and the widths range from the cent
floor for deep options to 20 cents. Each quote rests on the exchange until the next \omterm{def:hf:automated-options-market-making:mass}{mass quote} replaces it; between the underlying's move and that
replacement, all 170 are stale at once.

The exchange's answer is protection: a set of counters per market maker and per class (count of executions, contracts, notional, or the percentage
of quoted size executed) over a short window, and a rule that cancels all its quotes in the class when a counter reaches its limit.

\begin{dated}{2026-09}{One exchange's risk parameters}\label{dat:hf:automated-options-market-making:cboe}
Cboe's page on its options exchanges' risk management tools lists four parameter types: notional (contracts traded times premium), volume (contracts
traded), count (number of executions) and percentage of quote (the sum across all series of a root of the percentage of quoted contracts
executed), each over a period set anywhere from 100 milliseconds to a whole day; once a parameter is reached, no new trades are executed and any
trades in route are rejected.
\end{dated}

The chapter's sweep (\cref{lst:hf:automated-options-market-making:sweep}): the stock moves 1\%; 150 of the 170 quotes are now mispriced; an informed
trader hits them, most profitable first, one every five microseconds; the exchange pulls all quotes after $n$ executions within 100 milliseconds.
The loss grows almost linearly in $n$, about \$1\,000 an execution, because the most profitable targets are deep in-the-money options that move
dollar for dollar with the stock. On an ordinary day, customer orders arrive every five seconds on average, 70\% of them for one series and the rest
for spreads and strips of two, four and twelve legs executing within two milliseconds; each fill earns \$4. A threshold below twelve pulls the
quotes on legitimate twelve-leg orders and leaves the market maker out for half a second each time
(\cref{fig:hf:automated-options-market-making:protection}).

\begin{omfigure}
\begin{tikzpicture}
\begin{axis}[width=11cm,height=5.4cm,axis y line*=right,axis x line=none,xmode=log,log basis x=10,xmin=1.6,xmax=50,ymin=0,ymax=2600,
  ylabel={\small ordinary-day fills lost},tick label style={font=\footnotesize},table/col sep=comma]
\addplot[omDef,thick,dashed,mark=square*,mark size=1.3pt] table[x=n,y=lost]{figdata/hft/19-automated-options-market-making/protection.csv};
\label{plot:hf:protlost}
\end{axis}
\begin{axis}[width=11cm,height=5.4cm,axis y line*=left,xmode=log,log basis x=10,xmin=1.6,xmax=50,xtick={2,3,5,10,13,20,40},
  xticklabels={2,3,5,10,13,20,40},xlabel={\small protection threshold (executions in 100 ms)},ylabel={\small sweep loss (\$ thousand)},ymin=0,
  ymax=42,tick label style={font=\footnotesize},grid=major,grid style={gray!25},table/col sep=comma,
  legend style={font=\scriptsize,at={(0.5,-0.3)},anchor=north},legend columns=2]
\addplot[omThm,thick,mark=*,mark size=1.3pt] table[x=n,y=loss]{figdata/hft/19-automated-options-market-making/protection.csv};
\addlegendentry{loss to a 1\% sweep (left)}
\addlegendimage{/pgfplots/refstyle=plot:hf:protlost}\addlegendentry{fills lost on an ordinary day (right)}
\end{axis}
\end{tikzpicture}

\omcaption{The market maker's loss to an informed sweep after a 1\% move in the stock (left axis, thousand dollars) and the customer fills it loses
on an ordinary day while its quotes are pulled (right axis), against the exchange's count threshold; 170 series quoted ten contracts each; the
sweep hits one series every five microseconds; ordinary orders of one to twelve legs, 8\,235 fills a day without protection. Without protection the
sweep takes \$79\,370. Data: \texttt{hf\_options.protection\_table}.}\label{fig:hf:automated-options-market-making:protection}
\end{omfigure}

\section{Delta hedging at speed}

The sweep leaves the market maker short about 13\,000 shares of delta at a threshold of thirteen (it sold calls and bought puts after the stock
rose). The loss is already made: the options were sold at prices below the new value. The hedge's job is to stop the next move adding to it, and at
speed the hedge is sent the instant the fills arrive, in the underlying or in the most liquid correlated instrument, within the hedging band of
One Quant Book 5, chapter 26. Muravyev showed that inventory matters for option prices more than was thought: decomposing the price impact of
option trades into inventory risk and asymmetric information, he found both large and the inventory component larger, and order imbalances due to
inventory risk moving option prices five times more than earlier estimates.

\section{Per-strike and per-expiry risk}

\begin{definition}[Strike risk limit]\label{def:hf:automated-options-market-making:strike}
A \emph{strike risk limit}\index{strike risk limit} is a market maker's cap on the exposure (vega, gamma, or contracts) it will accumulate in one
strike or one expiry of a chain; when a fill breaches it, the market maker cancels or widens its quotes in that strike or expiry until the
exposure is reduced.
\end{definition}

Count-based protection counts executions, not risk. A volatility sweep, after implied volatility rises two points, hits 130 quotes of every
expiry and leaves the market maker short \$10\,600 of vega a volatility point for a loss of \$16\,500. A limit of \$1\,000 of vega per expiry,
enforced by the market maker's own system 50 microseconds after the breaching fill, cuts the loss to \$9\,900; enforced in 5 microseconds, to
\$6\,700. The limit's value depends on its speed: every microsecond of reaction is another stale quote hit.

\section{What the major firms have published}

Options market makers publish little about their quoting systems; what is public is the exchanges' rules and specifications (the protections of the
dated box, the mass-quote messages), the regulators' filings, and academic work on the data those firms generate. The chapter cites only those.

\section{Strategy files}

\begin{strategyfile}{Surface-fitted quoting across a listed chain}\label{strat:hf:automated-options-market-making:surface}
\sfield{Who pays you, and why} Options users (hedgers, directional traders, volatility traders) who need a price in any series at any time.
\sfield{Instruments and venues} A stock's or index's listed chain on every options exchange.
\sfield{Signal} Theoretical values from a live surface; the underlying's move through the surface's rule.
\sfield{Sizing and execution} \omterm{def:hf:automated-options-market-making:mass}{Mass quotes} at a width set by vega and the surface's uncertainty; refit the smile on a slower clock.
\sfield{Costs} Exchange fees, the width given up to competitors, adverse selection by traders who time executions.
\sfield{How it dies} Stale quotes after moves; competitors with faster refits.
\sfield{Horizon, capacity, infrastructure} Microseconds to the day; thousands of series; OPRA-scale market data.
\sfield{Backtest honestly} The surface as it was known at each quote; fills where timing traders would have taken them.
\sfield{Sources} Muravyev and Pearson (2020); One Quant Book 5, chapters 7, 8 and 26.
\end{strategyfile}

\begin{strategyfile}{Mass-quote making with quote protection}\label{strat:hf:automated-options-market-making:mass}
\sfield{Who pays you, and why} As above, with the exchange's protections limiting the loss to sweeps.
\sfield{Instruments and venues} Exchanges with mass-quote messages and risk parameters.
\sfield{Signal} The same theos; the protection's thresholds set from the largest legitimate order.
\sfield{Sizing and execution} Set the count threshold just above the largest multi-leg order (thirteen against twelve-leg orders in the model).
\sfield{Costs} Fills lost when the protection trips on legitimate flow (326 a day at a threshold of ten, three at thirteen).
\sfield{How it dies} Sweeps faster than the protection's first execution, and protections set too loose.
\sfield{Horizon, capacity, infrastructure} Milliseconds; one protection per class and exchange.
\sfield{Backtest honestly} The exchange's protection logic replayed on the day's executions, including multi-leg orders.
\sfield{Sources} The dated box; this chapter.
\end{strategyfile}

\begin{strategyfile}{Delta-hedged options book}\label{strat:hf:automated-options-market-making:delta}
\sfield{Who pays you, and why} Options customers through the spread; the hedge keeps the book's value from following the stock.
\sfield{Instruments and venues} The options book and the underlying (or its future).
\sfield{Signal} The book's delta after each fill.
\sfield{Sizing and execution} Hedge immediately beyond a band, in the cheapest instrument; aggregate across the chain first.
\sfield{Costs} The underlying's spread and fees; gamma between hedges.
\sfield{How it dies} Correlated fills that move the book's delta faster than the hedge can follow.
\sfield{Horizon, capacity, infrastructure} Microseconds to seconds; the underlying's order entry next to the options'.
\sfield{Backtest honestly} Hedge prices after the fill's timestamp plus the hedge latency.
\sfield{Sources} Muravyev (2016); One Quant Book 5, chapter 26.
\end{strategyfile}

\begin{strategyfile}{Event-aware options quoting}\label{strat:hf:automated-options-market-making:event}
\sfield{Who pays you, and why} Traders positioning for earnings or macro events, who pay for the implied move.
\sfield{Instruments and venues} Expiries spanning a known event.
\sfield{Signal} The event variance in the term structure (Book 5's event-variance arithmetic), against the stock's history of moves.
\sfield{Sizing and execution} Quote the event expiries on a surface with the event's variance; tighten vega limits before the event; pull or widen at
the announcement (chapter 18).
\sfield{Costs} Being wrong about the event's size; gamma across the announcement.
\sfield{How it dies} Surprises larger than the implied move.
\sfield{Horizon, capacity, infrastructure} Days to the event; seconds around it.
\sfield{Backtest honestly} The event calendar as known at the time.
\sfield{Sources} Chapter 18; One Quant Book 5, chapter 8.
\end{strategyfile}

\begin{strategyfile}{Vega-limited skew quoting}\label{strat:hf:automated-options-market-making:vega}
\sfield{Who pays you, and why} Buyers of protection (puts) and sellers of calls who keep the skew rich.
\sfield{Instruments and venues} Out-of-the-money puts and calls across expiries.
\sfield{Signal} The skew's level against its history and the surface fit.
\sfield{Sizing and execution} Quote the wings with per-strike and per-expiry vega limits enforced within microseconds.
\sfield{Costs} Tail risk in the wings; hedges in the underlying and in other expiries.
\sfield{How it dies} A volatility sweep: \$16\,500 unprotected in the model; the limits' speed decides the loss (\$9\,900 at 50 microseconds,
\$6\,700 at 5).
\sfield{Horizon, capacity, infrastructure} Days; fast limit enforcement.
\sfield{Backtest honestly} Limits applied with their real reaction time.
\sfield{Sources} This chapter.
\end{strategyfile}

\section{Tutorial: forty strikes in a millisecond}\label{tut:hf:automated-options-market-making:tut}

\begin{tutorial}
\textbf{Goal.} Quote a chain from a surface, then measure what an informed sweep takes against each protection threshold and what the threshold
costs on ordinary days. \textbf{End state:} \cref{fig:hf:automated-options-market-making:protection} and the numbers of the text.
\begin{tutsteps}
\item \textbf{The \omterm{def:hf:automated-options-market-making:mass}{mass quote}} from Book 5's theoretical values.
\omcode{code/firm/optquoter/firm_optquoter.py}{63}{74}{Theo and vega from \texttt{firm.optmm.theo}, a half-width in volatility points times vega, a cent floor.}\label{lst:hf:automated-options-market-making:quote}
\item \textbf{The sweep} against the exchange's count and the market maker's vega limits.
\omcode{code/firm/optquoter/firm_optquoter.py}{108}{126}{Stale quotes hit most profitable first; the exchange's count and the per-expiry vega limit stop it.}\label{lst:hf:automated-options-market-making:sweep}
\item \textbf{An ordinary day}: multi-leg customer orders against the same thresholds (\texttt{hf\_options.protection\_table}).
\item \textbf{The volatility sweep} with vega limits at two reaction speeds (\texttt{hf\_options.vol\_sweep}); the smile refit
(\texttt{hf\_options.refit}).
\end{tutsteps}
\textbf{What to change next.} Replace the count by a percentage-of-quote parameter; let the sweep arrive on several exchanges at once, each with
its own protection; run the chain on \texttt{firm.exchsim} with its mass-quote message.
\end{tutorial}

\begin{center}
\small
\begin{tabular}{@{}lrrrrrrr@{}}
\toprule
threshold (executions) & 2 & 5 & 10 & 13 & 20 & 40 & none\\
\midrule
loss to a 1\% sweep (\$) & 1\,980 & 4\,950 & 9\,899 & 12\,867 & 19\,773 & 38\,718 & 79\,370\\
ordinary-day fills lost & 2\,441 & 1\,052 & 326 & 3 & 0 & 0 & 0\\
edge lost a day (\$) & 9\,764 & 4\,208 & 1\,304 & 12 & 0 & 0 & 0\\
\bottomrule
\end{tabular}
\end{center}

\section{Build: the options quoter}\label{bld:hf:automated-options-market-making:optquoter}

\begin{build}
\bfield{Purpose} Quote a chain from a live surface in \omterm{def:hf:automated-options-market-making:mass}{mass quotes}, and measure protections and limits against sweeps and ordinary flow.
\bfield{Interface} \texttt{Chain(spot, strikes, expiries\_days)}, \texttt{surface(atm, term, skew, ref)}, \texttt{mass\_quote(chain, spot, surf,
half\_width\_vol, size)}, \texttt{refit\_slice(strikes, spot, tau, vols)}, \texttt{Protection(count, window\_ms)}, \texttt{sweep(quotes, jump,
chain, surf, protection, gap\_us, size, dvol, vega\_limit, react\_us)}, \texttt{normal\_day(n\_series, protection, seed)}, \texttt{bucket\_vega}.
Built on Book 5's \texttt{firm.optmm}, \texttt{firm.volpnl}, \texttt{firm.svi} and \texttt{firm.bs}.
\bfield{Rules} Prices per share, 100 shares a contract; the market maker is the counterparty of every sweep fill.
\bfield{Acceptance tests} \texttt{code/firm/optquoter/tests/}: quotes bracket theo with widths proportional to vega; the sweep stops at the count and
its loss equals the sum of the filled edges; a larger threshold never loses less; a vega limit reacting faster loses less; the ordinary day loses no
fills when the threshold exceeds every order's legs; the SVI refit recovers a slice within its noise.
\bfield{Stretch} Percentage-of-quote and notional parameters; several exchanges; the surface refit incrementally on each trade.
\end{build}

\begin{omsources}
\item D.~Muravyev, Order flow and expected option returns, \emph{Journal of Finance} 71(2), 2016, 673--708.
\item D.~Muravyev, N.~D.~Pearson, Options trading costs are lower than you think, \emph{Review of Financial Studies} 33(11), 2020, 4973--5014.
\item Cboe, US Options risk management tools (web page, accessed 2026).
\end{omsources}

\section{Exercises}

\begin{exercise}[$\star$]\label{exo:hf:automated-options-market-making:1}
The 30-day at-the-money call has a vega of 0.114 a volatility point. What is its width at half a point either side?
\end{exercise}

\begin{exercise}[$\star$]\label{exo:hf:automated-options-market-making:2}
A deep in-the-money call with delta one is quoted at \$20.01 offered, ten contracts. The stock rises \$1. What does a sweeper gain on it?
\end{exercise}

\begin{exercise}[$\star$]\label{exo:hf:automated-options-market-making:3}
Why must the count threshold exceed the largest legitimate multi-leg order?
\end{exercise}

\begin{exercise}[$\star\star$]\label{exo:hf:automated-options-market-making:4}
Why does the sweep loss grow almost linearly in the threshold?
\end{exercise}

\begin{exercise}[$\star\star$]\label{exo:hf:automated-options-market-making:5}
Why does a per-expiry vega limit enforced in 50 microseconds lose more than one enforced in 5?
\end{exercise}

\begin{exercise}[$\star\star$]\label{exo:hf:automated-options-market-making:6}
Muravyev and Pearson found that timed executions pay under 40\% of conventional spreads. What does that mean for a market maker's widths?
\end{exercise}

\begin{exercise}[$\star\star\star$]\label{exo:hf:automated-options-market-making:7}
\emph{Coding.} Rerun the sweep with a 2\% move. At a threshold of 13, what is the loss, and without protection?
\end{exercise}

\begin{exercise}[$\star\star\star$]\label{exo:hf:automated-options-market-making:8}
\emph{Find the flaw.} ``Our protection threshold is forty executions; we have never lost more than \$5\,000 to a sweep, so it is safe.''
\end{exercise}

\section{Problem: Forty Strikes in a Millisecond}

\begin{problem}[{Weekend problem --- forty strikes in a millisecond}]\label{pb:hf:automated-options-market-making:1}
An options market maker sets its protections and limits for a chain of 170 series.

\medskip\noindent\textbf{Part I --- Quotes.}
\begin{enumerate}
\item Define a \omterm{def:hf:automated-options-market-making:live}{live surface fit} and how its work is split between clocks.
\item How well does an SVI refit recover a noisy slice?
\item Define a \omterm{def:hf:automated-options-market-making:mass}{mass quote}.
\item Give the chain's widths.
\end{enumerate}

\medskip\noindent\textbf{Part II --- Protection.}
\begin{enumerate}[resume]
\item Summarise the dated box.
\item Describe the sweep and the ordinary day.
\item Give the sweep loss and the fills lost at thresholds of 5, 10, 13 and 20.
\item Why do the most profitable targets come first?
\end{enumerate}

\medskip\noindent\textbf{Part III --- Risk.}
\begin{enumerate}[resume]
\item What delta does the sweep leave, and what does the hedge do?
\item What did Muravyev find about inventory and option prices?
\item Define a \omterm{def:hf:automated-options-market-making:strike}{strike risk limit}.
\item Give the volatility sweep's losses with limits at two speeds.
\end{enumerate}

\medskip\noindent\textbf{Part IV --- The verdict.}
\begin{enumerate}[resume]
\item State the \emph{named result}: the loss on a planted informed sweep as a function of the quote-protection threshold, and the fills given up by
a tight threshold on normal days.
\item Which threshold would you choose, and why?
\item What would several exchanges, each with its own protection, change?
\item Why are the market maker's own limits needed besides the exchange's?
\item What did Muravyev and Pearson find about taking liquidity in options?
\item Which strategy file is most exposed to sweeps?
\item How would you set the vega limits per expiry?
\item In one sentence: what does quote protection trade off?
\end{enumerate}
\end{problem}

\section{Interview questions}

\begin{interviewq}[$\star$ \iqroles{trader}]\label{iq:hf:automated-options-market-making:1}
The stock jumps 3\% in a millisecond. What happens to your options quotes, in order?
\end{interviewq}

\begin{interviewq}[$\star\star$ \iqroles{developer}]\label{iq:hf:automated-options-market-making:2}
Design the path from an underlying tick to a \omterm{def:hf:automated-options-market-making:mass}{mass quote} of 5\,000 series. Where does the time go?
\end{interviewq}

\begin{interviewq}[$\star\star$ \iqroles{researcher}]\label{iq:hf:automated-options-market-making:3}
How would you decide when to refit a smile rather than move it with the underlying?
\end{interviewq}

\begin{interviewq}[$\star\star$ \iqroles{risk}]\label{iq:hf:automated-options-market-making:4}
Your protection tripped forty times today on one class. What do you investigate?
\end{interviewq}

\begin{interviewq}[$\star\star$ \iqroles{trader}]\label{iq:hf:automated-options-market-making:5}
Why might you quote puts wider than calls at the same vega?
\end{interviewq}

\begin{interviewq}[$\star\star\star$ \iqroles{researcher}]\label{iq:hf:automated-options-market-making:6}
Derive the expected loss of a sweep against a count threshold $n$ when targets arrive in order of decreasing edge $e_1\ge e_2\ge\cdots$.
\end{interviewq}
